% Compile: pdflatex
\documentclass[10pt,a4paper]{article}
\usepackage[top=1.8cm, bottom=1.8cm, left=2.2cm, right=2.2cm]{geometry}
\usepackage{tabularx}
\usepackage{graphicx}
\usepackage{setspace}
\usepackage{array}

\setstretch{1.08}
\setlength{\parindent}{0pt}
\setlength{\parskip}{2.5pt}

\begin{document}

\begin{center}
    {\fontsize{12}{14}\selectfont\textbf{SURAT PERNYATAAN KEPEMILIKAN CIPTAAN}}\\[2pt]
    {\fontsize{10}{12}\selectfont\textbf{DIREKTORAT JENDERAL KEKAYAAN INTELEKTUAL}}\\[1pt]
    {\fontsize{9}{11}\selectfont\textbf{KEMENTERIAN HUKUM DAN HAK ASASI MANUSIA REPUBLIK INDONESIA}}
\end{center}

\vspace{0.4em}
\hrule height 0.8pt
\vspace{0.6em}

Yang bertandatangan di bawah ini:

\vspace{0.2em}
\begin{tabularx}{\textwidth}{@{}p{0.5cm}p{3.5cm}X@{}}
1. & Nama & \textbf{Hafidz Rizqullah Prasetya} (Ketua Tim / Pencipta 1) \\
   & NIM / NIK & 24/535493/SV/24243 / 330221xxxxxxxxxx \\
   & Alamat / Kewarganegaraan & Bulaksumur, Caturtunggal, Sleman, D.I. Yogyakarta / WNI \\[3pt]
2. & Nama & \textbf{Matthew Hayunaji Priantara} (Anggota / Pencipta 2) \\
   & NIM / NIK & 24/536179/SV/24400 / 330221xxxxxxxxxx \\
   & Alamat / Kewarganegaraan & Bulaksumur, Caturtunggal, Sleman, D.I. Yogyakarta / WNI \\[3pt]
3. & Nama & \textbf{Akmal Manggala Putra} (Anggota / Pencipta 3) \\
   & NIM / NIK & 24/536182/SV/24402 / 330221xxxxxxxxxx \\
   & Alamat / Kewarganegaraan & Bulaksumur, Caturtunggal, Sleman, D.I. Yogyakarta / WNI \\
\end{tabularx}

\vspace{0.5em}
Dengan ini menyatakan bahwa Ciptaan yang diajukan permohonan pencatatannya adalah:
\vspace{0.2em}

\begin{tabularx}{\textwidth}{@{}p{4.2cm}X@{}}
\textbf{Jenis Ciptaan} & : \textbf{Program Komputer} \\
\textbf{Judul Ciptaan} & : \textbf{Verifin: Sistem Pendukung Keputusan Deteksi Lowongan Kerja Palsu Berbasis Bukti OSINT dan Explainable AI} \\
\textbf{Tanggal \& Tempat Diumumkan} & : 12 Agustus 2026 di Yogyakarta, Indonesia \\
\end{tabularx}

\vspace{0.4em}
Menyatakan dengan sesungguhnya bahwa:
\begin{enumerate}
    \setlength{\itemsep}{1pt}
    \setlength{\parskip}{0pt}
    \item Ciptaan tersebut di atas adalah \textbf{asli karya cipta kami sendiri}, bukan tiruan, jiplakan, ataupun saduran dari ciptaan orang lain tanpa izin sah;
    \item Ciptaan tersebut tidak bertentangan dengan peraturan perundang-undangan yang berlaku, ketertiban umum, kesusilaan, maupun norma keagamaan;
    \item Naskah dan perangkat lunak ini diikutsertakan dalam kompetisi \textbf{GEMASTIK XIX Tahun 2026} pada Cabang \textbf{Pengembangan Perangkat Lunak};
    \item Apabila di kemudian hari terbukti pernyataan ini tidak benar atau timbul sengketa dari pihak ketiga, kami bersedia menanggung segala akibat hukum yang timbul secara perdata maupun pidana.
\end{enumerate}

Demikian surat pernyataan ini dibuat dengan sebenar-benarnya untuk dipergunakan sebagaimana mestinya dalam permohonan pencatatan ciptaan di Ditjen Kekayaan Intelektual Kemenkumham RI.

\vspace{1em}

\begin{tabularx}{\textwidth}{@{}X c X@{}}
    Mengetahui, & & Yogyakarta, 12 Agustus 2026 \\[2pt]
    Dosen Pembimbing, & & Yang Menyatakan (Ketua Tim), \\[20pt]
    \textbf{Dr.Eng. Ir. Ganjar Alfian, S.T., M.Eng.} & \phantom{space} & \textbf{Hafidz Rizqullah Prasetya} \\
    NIP. 11198701202201101 & & NIM. 24/535493/SV/24243 \\
    & & \textit{\footnotesize (Meterai Rp10.000)} \\[10pt]
    \multicolumn{3}{c}{\textbf{Pencipta Lainnya:}} \\[6pt]
    \textbf{Matthew Hayunaji Priantara} & & \textbf{Akmal Manggala Putra} \\
    NIM. 24/536179/SV/24400 & & NIM. 24/536182/SV/24402 \\
\end{tabularx}

\end{document}
